\section{Discussion}

\subsection{What the ERA5 Signal Represents}
The results support a specific interpretation of ERA5 conditioning. ERA5 is not merely another correlated input channel that improves every setting equally. Its value grows as station outages become longer, precisely when local interpolation and station-history cues weaken. This is consistent with ERA5 acting as a physically coherent background state during multi-day sensor failures.

The variable-masking analysis further shows that the ERA5 contribution is not uniform across variables. Temperature, wind speed, and specific humidity produce the largest MAE increases when removed. Relative humidity is least important among the tested channels, which is consistent with the concern that ERA5 humidity can be biased in polar settings. This supports the decision to condition on ERA5 through missing-variable gates rather than replacing station observations wholesale.

\subsection{Why Complex Fusion Did Not Dominate}
The cross-attention and gated-injection experiments show that the specific fusion mechanism is not the dominant source of improvement. Gated injection and concat/no-cross fusion are statistically tied, while removing ERA5 causes a large degradation. This suggests that once aligned ERA5 information and block-missing training are available, a lightweight conditioning mechanism is sufficient for the present benchmark.

This finding changes the interpretation of the method. The contribution is not that attention-style fusion is intrinsically superior; instead, the contribution is identifying and validating the conditions under which external reanalysis becomes useful for Antarctic AWS imputation. Reporting the negative fusion result is important because it prevents over-claiming architectural novelty and makes the empirical conclusion more robust.

\subsection{Curriculum and Failure-Mode Matching}
The MCAR-on-block evaluation shows that ERA5 conditioning alone is insufficient. A model trained with MCAR masks performs poorly when evaluated on block-missing tests, even though it has access to ERA5. This indicates an interaction between the external physical prior and the missingness curriculum: the model must learn to use ERA5 under contiguous outage conditions, not only under randomly scattered missing points.

For practical Antarctic AWS recovery, this matters more than performance on generic random-missing benchmarks. Station failures often remove contiguous spans of one or more variables, so evaluation and training should match that failure mode.

The completed factorial matrix makes this conclusion quantitative. Block-missing curriculum has a larger average main effect than ERA5 conditioning (0.166 versus 0.105 MAE), while their effect-coded interaction is positive but small (0.017 MAE). This is methodologically useful beyond the present dataset: exogenous variables can be valuable, but their value depends on whether the training mask distribution teaches the model how to use them under the target failure mode.

\subsection{Practical Deployment Considerations}
The temporal-resolution analysis suggests that ECBIT does not require fragile hour-level ERA5 structure. Downsampling ERA5 to 6-hour resolution changes MAE by less than 0.001, while daily inputs are clearly too coarse. This supports practical deployments where only lower-frequency reanalysis products or delayed operational reanalysis streams are available, provided that sub-daily temporal structure is preserved.

The variable-importance results also suggest a prioritization rule for partial ERA5 availability. Temperature, wind speed, and specific humidity are the most valuable channels and account for about 84\% of the total single-channel masking degradation; pressure and relative humidity contribute less in the current benchmark. If a deployment pipeline must reduce ERA5 variables, the top three channels should be preserved first.

The held-out station results further suggest that deployment should include station-level validation. Mean MAE ranges from 0.173 to 0.350 across the five held-out stations, so high-elevation, coastal, or locally complex stations should be checked before operational use rather than assumed to follow the aggregate benchmark average.

\subsection{Station Heterogeneity}
Held-out station results vary by roughly a factor of two, from Nico to Mount Sidley. This heterogeneity likely reflects a combination of local geography, elevation, coastal versus interior climate regimes, and station-specific ERA5 representativeness. Mount Sidley has the largest held-out MAE. In the selected metadata it is also the highest held-out station, at 2123 m, and lies in West Antarctic complex terrain where coarse reanalysis orography and boundary-layer structure can be mismatched to the local AWS footprint. Zhongshan has the second-largest MAE despite its low elevation, plausibly because its coastal setting is affected by sea-ice, ocean, and local wind interactions that are not fully resolved by point-sampled ERA5. Future work should therefore consider station-adaptive conditioning, metadata-aware calibration, or few-shot adaptation for stations with large ERA5 bias.

\subsection{Limitations}
The benchmark uses artificial masks on originally observed positions, which provides controlled ground truth but does not directly evaluate recovery inside historical gaps where no AWS observation exists. ERA5 is aligned at station points and shares the selected variable set; other reanalysis products, spatial interpolation strategies, or station metadata may change the optimal conditioning strategy. Finally, the completed result matrix prioritizes a controlled 2$\times$2 curriculum-by-ERA5 factorial design over a broad third-party imputation suite. BRITS and SAITS wrappers are implemented, but complete block-curriculum retraining and evaluation for those methods is left for future benchmarking; the current study isolates why ERA5 and block-missing curriculum matter before expanding the model zoo.
